\section{The supporting map and its relative commutant}
\label{sec:commutant}

We retain the notation of the preceding section.  In particular,
\[
 R=A\otimes\Z_1\otimes\Z_2\otimes\Z_3,
 \qquad B=R_\omega,
\]
where $A$ is nonzero, separable, unital and nuclear.  The maps
$\pi,\theta\colon A\to B$ have ranges in the first subsystem
$(A\otimes\Z_1)_\omega$, with $\pi$ the constant first-factor
homomorphism.  The positive contractions $e,q=1-e$ satisfy
\eqref{setup:common-equations}.  The second-leg contraction $P$ defines
\[
 \Psi(a)=P(q-1/2)_+\pi(a),\qquad a\in A.
\]
We shall use both the cut comparisons \eqref{trace:small-part-comparison} and
the arbitrarily long finite families of second-leg elements in
\eqref{trace:ultrapower-copies}.  The constant third leg will remain available
for the functional-calculus argument.

\subsection{A supporting map preserving the relevant ideals}

Put $h=\theta(1_A)$.  An abstract supporting homomorphism for an
order-zero map need not take values in its codomain.  We therefore seek
a c.p.c. order-zero map $\widehat\theta\colon A\to B$ with three
roles.  Its factorization $h\widehat\theta=\theta$ recovers the model
and gives $h\widehat\theta(1_A)=h$, the support identity for $h$ in
the second matrix corner below.  The equality
$q\widehat\theta=q\pi$ preserves the common supported part.
Finally, pointwise ideal containment of $e\widehat\theta(a')$ and
$e\pi(a')$ in $\Ideal_B(\Psi(a'))$ supplies the hypothesis for
Gabe's theorem, Theorem~\ref{prelim:domination}, to approximate the
error maps by finite sums of compressions of $\Psi$.

\begin{proposition}[A supporting map preserving ideals]\label{support:map}
There is a c.p.c. order-zero map $\widehat\theta\colon A\to B$
whose range commutes with the constant copy of $\Z_3$ and such that
\begin{align}
 h\widehat\theta(a)=\widehat\theta(a)h&=\theta(a),
       \qquad a\in A,\label{support:factorization}\\
 [e,\widehat\theta(a)]&=0,
 \qquad q\widehat\theta(a)=q\pi(a),\qquad a\in A,
       \label{support:common}\\
 e\widehat\theta(a'),\ e\pi(a')&\in\Ideal_B(\Psi(a')),
       \qquad a'\in A_+.
       \label{support:ideals}
\end{align}
\end{proposition}

\begin{proof}
Consider the order-zero root maps $\theta^{1/l}$, for $l\geq1$.
They commute with $e$ and the constant third leg.  Indeed, each value
$\theta^{1/l}(a)$ is a norm limit of sandwiches of $\theta(a)$ by
bounded continuous functions of $h$.  One obtains these sandwiches
by truncating $t^{(1/l-1)/2}$ from above; the resulting scalar weights
converge uniformly to $t^{1/l}$ after multiplication by $t$.
The same observation gives
\begin{equation}\label{support:root-comparison}
 e\theta^{1/l}(x_j)\precsim e\theta(x_j).
\end{equation}
Here $e$ commutes with the values involved, so the sandwiches also
approximate the positive element on the left.

Taking $a=1_A$ in \eqref{setup:common-equations} gives $qh=q$.  Continuous
functional calculus therefore gives $qh^{1/l}=q$, and the sandwich
formula together with \eqref{setup:common-equations} yields
\[
 q\theta^{1/l}(a)=q\pi(a),\qquad a\in A.
\]
Moreover,
\[
 h\theta^{1/l}(a)=\theta^{1/l}(a)h\longrightarrow\theta(a)
 \quad(l\longrightarrow\infty),
\]
by the uniform convergence of $t^{1+1/l}$ to $t$ on $[0,1]$.

For every fixed $j$, combining \eqref{support:root-comparison} with
\eqref{trace:small-part-comparison} gives
\[
 e\theta^{1/l}(x_j)\precsim(\Psi(a_j)-c_j)_+,
 \qquad c_j=\gamma_j/64>0.
\]
The bounded-witness estimate of Lemma~\ref{prelim:bounded-witness}, with the
spare cut $c_j$, now shows that this positive element can be
approximated arbitrarily closely by $r_j^*\Psi(a_j)r_j$, with
\begin{equation}\label{support:witness-bound}
 \|r_j\|^2\leq 2/c_j.
\end{equation}
The bound depends on $j$ but, crucially, not on $l$.

We apply the countable diagonalization principle of
Lemma~\ref{prelim:epsilon-test}.  To make the boundedness issue explicit,
lift each root map to a sequence of c.p.c. maps into $R$ by the
Choi--Effros lifting theorem.  Keep fixed representatives of
$e,q,h,\theta,\Psi$ and $\pi$.  The tests impose the factorization
and common-part identities on countable dense subsets of $A$,
commutation with $e$ and a dense subset of the constant third leg,
and order zero on a countable dense family of orthogonal pairs of
positive contractions.  For each $j$ include a witness variable with
a fixed representative bound strictly larger than
$\sqrt{2/c_j}$, and impose
\[
 e\widehat\theta(x_j)=r_j^*\Psi(a_j)r_j.
\]
Every finite list is approximately solvable: take $l$ sufficiently
large for the factorization tests, and then use the witnesses above
with sufficiently small errors.  Their fixed bounds allow the
ultrapower tests to be realized by bounded coordinate variables.
The diagonalization produces a c.p.c. map satisfying all the tests.
Continuity extends the identities and commutations to all the
specified elements; the dense orthogonality tests make the map
order zero.  This proves \eqref{support:factorization} and
\eqref{support:common}, as well as the witness identities for all $j$.

It remains to extend ideal containment from the listed cuts in
\eqref{model:tests} to an arbitrary positive contraction $a'\in A$.
Such an $a'$ can be
approximated arbitrarily closely by $x_j$ with
$a_j\in\Ideal_A(a')$.  To see this, choose $d$ in the fixed dense
positive family close to $a'$, and a small positive rational
$r>\|d-a'\|$.  The image of $d$ in $A/\Ideal_A(a')$ has norm less
than $r$, so $(d-r)_+\in\Ideal_A(a')$.  This cut belongs to the
chosen family; take it as $a_j$ and choose the accompanying rational
$\gamma_j$ small.  Then $x_j=(a_j-\gamma_j)_+$ is as close to $a'$
as desired.

The map $\Psi$ is a positive weight times the homomorphism $\pi$,
and the weight commutes with $\pi(A)$.  Consequently
\[
 a_j\in\Ideal_A(a')
 \quad\Longrightarrow\quad
 \Psi(a_j)\in\Ideal_B(\Psi(a')):
\]
apply $\Psi$ to finite sums of products through $a'$ and pass to
norm limits.  The witness identity places
$e\widehat\theta(x_j)$ in this ideal.  The comparison
\eqref{trace:small-part-comparison} similarly places $e\pi(x_j)$ there.
Taking norm limits, and then scaling, proves
\eqref{support:ideals} for every $a'\in A_+$.
\end{proof}

\subsection{Comparison inside the supported relative commutant}

Define the c.p.c. order-zero map
\[
 \rho=\operatorname{diag}(\rho_0,\rho_1)\colon A\longrightarrow M_2(B),
 \qquad \rho_0=\pi,\quad\rho_1=\widehat\theta,
\]
and its supported relative commutant
\begin{equation}\label{comm:definition}
 C=M_2(B)\cap\rho(A)'\cap\{1-\rho(1_A)\}^{\perp}.
\end{equation}
The last condition means annihilation on both sides.  A matrix in $C$
supported in position $(0,1)$ has an entry that intertwines $\rho_1$
with $\rho_0$ and satisfies both support conditions.  The matrix's two
modulus squares lie in the diagonal corners of $C$ and have the same Cuntz
class, allowing comparison of positive elements in the two corners
inside one semigroup.  For $b\in B$,
write $b^{(i)}$ for the matrix with entry $b$ in position $(i,i)$ and
zeros elsewhere, using indices $i=0,1$.  Put
\[
 C_i=\{b\in B:b^{(i)}\in C\},\qquad
 Q_i=q^{(i)},\qquad H=1_B^{(0)}.
\]
The common-part identities show that $Q_i,H\in C$.  They also imply
that the matrix with $q^{1/2}$ in position $(0,1)$ belongs to $C$:
positive-support functional calculus passes the intertwining and
support identities from $q$ to $q^{1/2}$.  Its two modulus squares
are $Q_1$ and $Q_0$.  Thus
\begin{equation}\label{comm:common-equivalence}
 [Q_0]=[Q_1]\quad\text{in }W(C).
\end{equation}

\begin{proposition}[Packing the error part]\label{comm:packing}
For $i\in\{0,1\}$ and every positive contraction $c\in C_i$,
\begin{align}
 m[(c^{1/2}ec^{1/2})^{(i)}]&\leq[Q_0],
       &&m\in\N,\label{comm:small-part}\\
 [c^{(i)}]&\leq2[Q_0]
       &&\text{in }W(C).\label{comm:two-common}
\end{align}
\end{proposition}

\begin{proof}
Fix $m\geq1$ and set $t=c^{1/2}ec^{1/2}$.  This positive
contraction commutes with $\rho_i(A)$ and satisfies
$t\rho_i(1_A)=t$.  The order-zero product identity therefore gives
\begin{equation}\label{comm:supported-product}
 t\rho_i(a)\rho_i(b)=t\rho_i(ab),\qquad a,b\in A.
\end{equation}
Apply Theorem~\ref{prelim:domination} to
\[
 \Phi(a)=t\rho_i(a),\qquad
 \Psi(a)=P(q-1/2)_+\pi(a).
\]
Both maps are completely positive and nuclear, since their domain
is nuclear, and $\Psi$ is order zero.  For $a\in A_+$,
\[
 \Phi(a)=c^{1/2}e\rho_i(a)c^{1/2}
       \in\Ideal_B(\Psi(a))
\]
by Proposition~\ref{support:map}.  Hence finite sums
$\sum_{l=1}^L y_l^*\Psi(\,\cdot\,)y_l$ approximate $\Phi$
pointwise in norm.

Write $h_q=(q-1/2)_+$.  For one such approximation, take $mL$
second-leg elements $v_{sl}$ from \eqref{trace:ultrapower-copies} and set
\[
 r_s=\sum_{l=1}^L v_{sl}h_q^{1/2}y_l,
 \qquad 1\leq s\leq m.
\]
These second-leg elements commute with $q$ and $\pi(A)$, and satisfy
$v_{sl}^*v_{s'l'}=\delta_{ss'}\delta_{ll'}P$.  Expanding the products
therefore gives
\begin{equation}\label{comm:packed-approximation}
 r_s^*\pi(a)r_{s'}=
 \delta_{ss'}\sum_{l=1}^L y_l^*\Psi(a)y_l,
 \qquad a\in A.
\end{equation}
For $f(u)=\min\{1,2u\}$ on $[0,1]$ one also has
$f(q)r_s=r_s$, since $f(q)h_q^{1/2}=h_q^{1/2}$.
Include $1_A$ in the approximation set.  As
$\Phi(1_A)=t$ and $\|t\|\leq1$, sufficiently accurate
approximations in \eqref{comm:packed-approximation} have
$\|r_s\|^2\leq2$.  Thus, for this fixed $m$, there is a fixed
bound on every packed variable, independent of $L$ and of the
approximation accuracy.

Using bounded coordinate lifts and the diagonalization principle of
Lemma~\ref{prelim:epsilon-test} on a countable dense subset of $A$, we
obtain $r_1,\ldots,r_m\in B$ with
\begin{equation}\label{comm:packed-exact}
 f(q)r_s=r_s,\qquad
 r_s^*\pi(a)r_{s'}=\delta_{ss'}t\rho_i(a),
 \qquad a\in A.
\end{equation}
These elements are intertwiners.  Indeed, using
\eqref{comm:supported-product} and
\eqref{comm:packed-exact}, the right modulus square of
$\pi(a)r_s-r_s\rho_i(a)$ is
\[
 t\rho_i(a^*a)-t\rho_i(a^*)\rho_i(a)=0.
\]
Consequently $\pi(a)r_s=r_s\rho_i(a)$ for all $a\in A$.
Taking $a=1_A$ gives $r_s\rho_i(1_A)=r_s$ as well.
The matrix $R_s$ with $r_s$ in position $(0,i)$ thus belongs to
$C$, including its support condition.

For the row $V=(R_1,\ldots,R_m)$ over $C$, the positive matrix
$V^*V$ has $m$ diagonal entries equal to $t^{(i)}$ and all
other entries zero.  Meanwhile $VV^*\in\her_C(Q_0)$:
\eqref{comm:packed-exact} gives
$f(Q_0)VV^*f(Q_0)=VV^*$, and $f(0)=0$.
Equality of the Cuntz classes of the two modulus squares now proves
\eqref{comm:small-part}.

Finally,
\[
 c=c^{1/2}ec^{1/2}+c^{1/2}qc^{1/2}.
\]
The second summand, placed in corner $i$, is Cuntz below $Q_i$
inside $C$.  Apply \eqref{comm:small-part} with $m=1$, the
subadditivity of Cuntz classes of positive sums, and
\eqref{comm:common-equivalence}.  This proves
\eqref{comm:two-common}.
\end{proof}

\begin{corollary}\label{comm:fullness}
The elements $Q_0,Q_1$ and $H$ are full in $C$, and $C$ is full
in $M_2(B)$.
\end{corollary}

\begin{proof}
Every positive element of $M_2(B)$ is bounded by twice its diagonal
part.  For an element of $C_+$, each diagonal corner also lies in
$C$, because $\rho$ and its support condition are diagonal.
After scaling, \eqref{comm:two-common} shows that these corners
belong to the ideal generated by $Q_0$.  Hence so does the original
positive element.  Thus $Q_0$ is full; the same holds for $Q_1$ by
\eqref{comm:common-equivalence}, and for $H$ because $Q_0\leq H$.
The corner projection $H$ is full in the ambient matrix algebra
$M_2(B)$, so the latter also contains $C$ as a full subalgebra.
\end{proof}

\begin{proposition}\label{comm:unperforated}
The ordered semigroup $W(C)$ is almost unperforated.
\end{proposition}

\begin{proof}
Suppose $(N+1)x\leq Ny$ in $W(C)$, with $N\geq1$.
Choose a separable $C^*$-subalgebra $C_*\subset C$ containing the
matrix entries of representatives of $x,y$ and of countably many
comparison witnesses.  The inequality then already holds over
$C_*$.

There is a unital copy of $\Z$ in $M_2(B)$ commuting with both
$C_*$ and $\rho(A)$.  Here is the coordinate argument, which also
allows the representatives of $C_*$ to vary with the coordinate.
Choose isomorphisms
$\sigma_l\colon\Z\otimes\Z\to\Z$ with
$\sigma_l(d\otimes1)\to d$, and put
$\nu_l(z)=\sigma_l(1\otimes z)$.  For fixed $d\in\Z$,
\[
 \|[\nu_l(z),d]\|
 \leq 2\|z\|\,\|d-\sigma_l(d\otimes1)\|.
\]
Thus the commutators tend to zero uniformly for $z$ in the unit
ball.  Place these embeddings diagonally on the third tensor leg
of $M_2(R)$.  Tensor approximation gives the same uniform
commutation property with any fixed element of $M_2(R)$.
At coordinate $n$, choose $l$ sufficiently large for the entries
of representatives of the first $n$ elements of countable dense
commutation data, with tolerance tending to zero.  The resulting
coordinate embeddings induce a unital embedding
$\nu\colon\Z\to M_2(B)$ commuting with $C_*$ and $\rho(A)$.
It is isometric because each coordinate embedding is isometric.

Multiplication gives a homomorphism
\[
 C_*\otimes\Z\longrightarrow C,\qquad c\otimes z\longmapsto c\nu(z).
\]
Commutation gives a homomorphism on the maximal tensor product;
nuclearity of $\Z$ identifies it with the spatial tensor product.
The range satisfies the support condition because $c$ does and
$\nu(z)$ commutes with $\rho(1_A)$.  This homomorphism restricts
to the original inclusion on $C_*\otimes1$.

The algebra $C_*\otimes\Z$ is $\Z$-stable.  Its Cuntz semigroup
is therefore almost unperforated by Theorem~\ref{prelim:comparison}.  Tensoring the original inequality with
$1_\Z$ gives comparison of the images of $x$ and $y$ there.
Mapping the comparison back to $C$ yields $x\leq y$.

Here the reindexed copy $\nu(\Z)$ is used for this inequality alone.
The constant third-leg copy retained above is the separate source of
the weight $k$ below; no commutation between these two copies is
required.
\end{proof}

We next use the full projection $H$, the corner divisibility supplied
by Lemma~\ref{supplies:copies}, and almost unperforation of $W(C)$
to verify the structural hereditary hypothesis of the BBSTWW
criterion recalled in Lemma~\ref{prelim:unitary}.

\begin{proposition}\label{conjugacy:hereditary}
Let $D\subset C$ be a full hereditary subalgebra.  Suppose that
$y\in D$ and that totally full positive elements $e_l,e_r\in D$
satisfy
\[
 e_l y=ye_r=0.
\]
Then there is a full element $s\in D$ with $sy=ys=0$.
\end{proposition}

\begin{proof}
Fullness of $e_l,e_r$ in $D$ and fullness of $D$ in $C$ imply
that both elements are full in $C$.  Since $H$ is a projection,
there is a common finite integer $N\geq1$ such that
\begin{equation}\label{conjugacy:projection-domination}
 [H]\leq N[e_l],\qquad [H]\leq N[e_r]
 \quad\text{in }W(C).
\end{equation}
Use the positive contraction $D_{N+1}$ from the order-zero supplies
on a constant $\Z$ leg, and put
$p_0=(D_{N+1})^{(0)}$.  It belongs to $C$, since the first-corner
map is $\pi$, and it has norm one.  Simplicity of the constant
leg gives fullness of $p_0$ in the corner with unit $H$; because
$H$ is full in $C$, so is $p_0$.

The $N+1$ accompanying elements from Lemma~\ref{supplies:copies},
placed in this corner, form a row over $C$ whose right modulus
square is the diagonal sum of $N+1$ copies of $p_0$, while its
left modulus square belongs to $HCH$.  Hence
\[
 (N+1)[p_0]\leq[H].
\]
Combining this with \eqref{conjugacy:projection-domination} and
Proposition~\ref{comm:unperforated} gives
$p_0\precsim e_l,e_r$ in $C$.  We may therefore choose
$w_l,w_r\in C$ so that
\[
 A_l=w_l^*e_lw_l,\qquad A_r=w_r^*e_rw_r
\]
are as close to $p_0$ as prescribed.

We specify the accuracy before choosing these witnesses.
Fullness of $p_0^2$ in the constant simple leg gives fixed
$z_1,\ldots,z_L\in HCH$ with
\[
 \sum_{j=1}^L z_j^*p_0^2z_j=H.
\]
Choose $\varepsilon>0$ such that
\begin{equation}\label{conjugacy:fullness-tolerance}
 (2\varepsilon+\varepsilon^2)
       \sum_{j=1}^L\|z_j\|^2<1,
\end{equation}
and choose $w_l,w_r$ with $\|A_l-p_0\|,\|A_r-p_0\|<\varepsilon$.
Set
\[
 s=e_r^{1/2}w_rw_l^*e_l^{1/2}.
\]
Heredity gives $s\in D$, since $DCD\subset D$.
Functional calculus applied to the annihilation identities gives
$e_l^{1/2}y=ye_r^{1/2}=0$, so $sy=ys=0$.

Let $J=\Ideal_C(s)$.  It contains
\[
 w_r^*e_r^{1/2}s\,e_l^{1/2}w_l=A_rA_l,
 \qquad
 \|A_rA_l-p_0^2\|\leq2\varepsilon+\varepsilon^2.
\]
Thus the element $b=\sum_j z_j^*A_rA_lz_j\in J$ satisfies
$\|H-b\|<1$ by \eqref{conjugacy:fullness-tolerance}.
The image of $H$ in $C/J$ is a projection of norm less than one,
and is therefore zero.  Since $H$ is full in $C$, this forces
$J=C$.

It remains to prove that $s$ is full inside $D$.  Let
$I=\Ideal_D(s)$.  For any ideal $I$ of a hereditary
subalgebra $D\subset C$, one has
\[
 D\,\Ideal_C(I)\,D\subset I.
\]
Indeed, the ideal in the middle is the closed linear span of
products through elements of $I$.  By linear decomposition it is
enough to use a positive middle element $i\in I_+$, and then
\[
 dci c'd'
   =(dc\,i^{1/3})\,i^{1/3}\,(i^{1/3}c'd')\in I:
\]
the two outer factors lie in $D$ by heredity.
Here $\Ideal_C(I)$ contains $J=C$, so the displayed inclusion
and an approximate identity of $D$ give $D\subset I$.
Hence $I=D$, proving that $s$ is full in $D$.
\end{proof}

\subsection{Full-spectrum weights and unitary equivalence}

The structural conditions on $C$ in Lemma~\ref{prelim:unitary} are now
verified.  We turn to the remaining input: totally full contractions
with matching Cuntz classes under positive functional calculus.

Fix a positive contraction $k$ of spectrum $[0,1]$ in the retained
constant third leg.  Such a $k$ is obtained from the scalar interval
coordinate in a prime dimension-drop algebra embedded in $\Z_3$.
When necessary, a third-leg element also denotes its diagonal
copy in $M_2(B)$.  Set
\begin{equation}\label{conjugacy:weights}
 h_0=k^{(0)},\qquad h_1=(kh)^{(1)}.
\end{equation}
These are positive contractions in $C$.  In the second corner,
the support condition follows from
$h\widehat\theta(1_A)=h$; commutation follows from
Proposition~\ref{support:map}.

\begin{proposition}\label{conjugacy:functions}
The elements $h_0,h_1$ are totally full in $C$.  For every
$g\in C_0((0,1])_+$, the elements $g(h_0)$ and $g(h_1)$ are
Cuntz equivalent in $C$.
\end{proposition}

\begin{proof}
The case $g=0$ is immediate.  Otherwise rescale so that
$\|g\|\leq1$, and put
\[
 G_i=g(h_i),\qquad T_i=Q_iG_i,\qquad E_i=G_i-T_i=e^{(i)}G_i.
\]
All these elements are positive and belong to $C$.  The factors
commute because $e$ commutes with $h$ and is represented in the
first subsystem.  Since $qh=q$, functional calculus gives
$T_i=Q_i g(k)$.  The matrix with
$q^{1/2}g(k)^{1/2}$ in position $(0,1)$ belongs to $C$ and has
modulus squares $T_1,T_0$.  Consequently
\begin{equation}\label{conjugacy:common-functions}
 [T_0]=[T_1]\quad\text{in }W(C).
\end{equation}

The positive element $g(k)$ is nonzero in the simple unital
constant third leg.  Choose finitely many elements
$z_1,\ldots,z_K$ in that leg with
\[
 \sum_{l=1}^K z_l^*g(k)z_l=1.
\]
For $\delta>0$ let
\[
 b_\delta(u)=
 \begin{cases}
 ((u-\delta)_+/u)^{1/2},&u>0,\\
 0,&u=0.
 \end{cases}
\]
The diagonalized $z_l$ commute with $Q_i$ and $\rho(A)$ and
multiply $C$ into itself.  In particular
$z_l b_\delta(Q_i)\in C$, and
\[
 (Q_i-\delta)_+
 =\sum_{l=1}^K
   (z_l b_\delta(Q_i))^*T_i(z_l b_\delta(Q_i)).
\]
Taking all cuts proves
\begin{equation}\label{conjugacy:fullness-bound}
 [Q_i]\leq K[T_i]\leq K[G_i]\quad\text{in }W(C).
\end{equation}
Since $Q_i$ is full, this shows that $g(h_i)$ is full for every
nonzero $g$.  In particular the two elements in
\eqref{conjugacy:weights} are totally full; the same argument
also ensures that every such functional calculus is nonzero.

Let $d$ be any dimension function on $W(C)$ normalized at either
$[G_0]$ or $[G_1]$, if one exists.  Equations
\eqref{comm:two-common} and
\eqref{conjugacy:fullness-bound}, together with
\eqref{comm:common-equivalence}, imply
\[
 0<d([Q_0])=d([Q_1])<\infty.
\]
Apply \eqref{comm:small-part} with $c^{(i)}=G_i$.
Since $G_i$ commutes with $e^{(i)}$, it follows that
$m[E_i]\leq[Q_0]$ for every $m$, and hence $d([E_i])=0$.
The inequalities
\[
 [T_i]\leq[G_i]\leq[T_i]+[E_i]
\]
and \eqref{conjugacy:common-functions} now give
\begin{equation}\label{conjugacy:dimension-equality}
 d([G_0])=d([G_1])=1.
\end{equation}

Fix $\eta>0$.  By continuity, $g(0)=0$, and $g\ne0$, there is
a nonzero continuous function $g'$ with $0\leq g'\leq g$ and
support contained in the open set where $0<g<\eta$.
For the dimension functions under consideration,
\eqref{conjugacy:fullness-bound} applied to $g'$ gives
$d([g'(h_i)])>0$.  This element is orthogonal to
$(G_i-\eta)_+$, and their sum is bounded by $G_i$.
Therefore \eqref{conjugacy:dimension-equality} gives
\[
 d([(G_i-\eta)_+])+d([g'(h_i)])\leq1,
 \qquad
 d([(G_i-\eta)_+])<1=d([G_{1-i}]).
\]
This holds for every dimension function normalized at
$[G_{1-i}]$.  The target $G_{1-i}$ is full, so the ideal condition
for the comparison criterion in Theorem~\ref{prelim:comparison} holds.
Proposition~\ref{comm:unperforated} and that criterion yield
$(G_i-\eta)_+\precsim G_{1-i}$.  This use of the criterion
also permits an empty family of normalized dimension functions.
Taking all positive cuts, and then exchanging $i=0,1$, proves
the required Cuntz equivalence.
\end{proof}

\begin{proposition}[Weighted unitary equivalence]\label{conjugacy:weighted}
Put $k^{[0]}=k$ and $k^{[1]}=1-k$, in the constant third leg.
There are unitaries $U_0,U_1\in M_2(B)$ such that
\begin{equation}\label{conjugacy:weighted-equation}
 U_i\bigl(k^{[i]}\theta(a)\bigr)^{(1)}U_i^*
   =\bigl(k^{[i]}\pi(a)\bigr)^{(0)},
 \qquad a\in A,\quad i=0,1.
\end{equation}
They can be chosen to commute with $\rho(A)$.
\end{proposition}

\begin{proof}
Apply the relative-commutant unitary-equivalence criterion of
Lemma~\ref{prelim:unitary}, namely \cite[Lemma~5.3]{BBSTWW}, to
\eqref{comm:definition}.  Its coordinate algebras are $M_2(R)$,
which are separable, unital and $\Z$-stable, and its domain is the
separable unital algebra $A$.  The map $\rho$ is c.p.c. order
zero.  Corollary~\ref{comm:fullness} verifies fullness of $C$ in
the ambient ultrapower; Proposition~\ref{conjugacy:hereditary}
verifies the required condition in every full hereditary
subalgebra.  Proposition~\ref{conjugacy:functions} supplies total
fullness and all the functional-calculus comparisons.  Thus there
is a unitary in the unitization of $C$ conjugating $h_1$ to $h_0$.

Realizing this unitary in the ambient algebra as explained after
Lemma~\ref{prelim:unitary} gives $U_0\in M_2(B)$ commuting with
$\rho(A)$. Multiplying
$U_0h_1U_0^*=h_0$ by $\rho(a)$ and using
\eqref{support:factorization} gives
\[
 U_0(k\theta(a))^{(1)}U_0^*=(k\pi(a))^{(0)}.
\]
The element $1-k$ also has spectrum $[0,1]$ in the constant
third leg.  Repeating the same argument with this element, with
all the previously constructed maps unchanged, gives $U_1$ and
the other identity in \eqref{conjugacy:weighted-equation}.
\end{proof}

Complementary weights are part of the two-coloured approximation
framework developed in \cite{SWW,BBSTWW,CETWW}. Here the required
ideal-sensitive comparisons have been proved above. The two identities
in \eqref{conjugacy:weighted-equation} add to the first-corner inclusion
because $k+(1-k)=1$.  The next section
passes from these ultrapower identities to finite-dimensional
approximations with exactly two order-zero return maps.
